\documentclass[11pt]{article}
\usepackage[a4paper,margin=25mm]{geometry}
\usepackage{amsmath,amssymb}
\usepackage[hidelinks]{hyperref}
\setlength{\emergencystretch}{3em}
\title{Conjecture 00000001674: no edge-count bound depending only on the excluded grid}
\author{lizaixi01, with OpenAI Codex assistance}
\date{23 October 2026}
\begin{document}
\maketitle
\section{The literal upper-bound premise is false}
The source explicitly defines its grid theorem as an upper bound
on the \emph{number of edges} of graphs excluding a $k\times k$
grid minor, and calls $2^{O(k\log k)}$ the optimal bound.
There is no finite edge-count upper bound depending only on $k$,
even for connected trees. For every $k\geq2$ and every finite
bound $B$, the actual star with $B+1$ leaves has $B+1$ edges
and excludes the $k\times k$ grid as a minor.

This is a statement about the literal edge count in the printed
definition. The genuine excluded-grid theorem controls
\emph{treewidth}, and is not being refuted. Nor do we deny a
bound with additional dependence on the number of vertices or a
properly normalized density. The source's separate existence
lower-bound clause is not independently contradicted; the
necessary finite upper-bound premise already fails.

\section{Actual stars and the standard minor model}
The host graph is Mathlib's actual \texttt{starGraph}, on
$\operatorname{Option}(\operatorname{Fin}n)$ with center
$r=\operatorname{none}$ and $n$ leaves. Its actual adjacency is
\[
 x\sim y\quad\Longleftrightarrow\quad
 x\ne y\ \land\ (x=r\ \lor\ y=r).
\]
Lean uses the actual theorem that this graph is a tree and the
actual finite-tree identity
$|E|+1=|V|$. Since $|V|=n+1$, its true
\texttt{edgeFinset.card} equals $n$. No edge count is assigned.

A minor model for a graph $A$ in a host $B$ consists of a branch
set $B_v$ for every vertex $v$ of $A$, with all standard conditions:
each $B_v$ is nonempty; its actual induced subgraph is connected;
distinct branch sets are disjoint; and for every edge $uv$ of $A$
there exist $x\in B_u$ and $y\in B_v$ with an actual host edge
$xy$. This is the standard single-cross-edge branch-set
definition of a minor. The connectedness and nonemptiness
conditions are retained even though the stronger obstruction
below needs only the last two conditions.

\section{Every grid of size at least two is excluded}
The actual grid has vertices
$\operatorname{Fin}k\times\operatorname{Fin}k$ and horizontal
or vertical adjacency exactly when the other coordinate is equal
and the changing natural coordinate values are consecutive.
For every $k\geq2$, it contains the two edges
\[
 (0,0)(0,1),\qquad (1,0)(1,1),
\]
whose endpoints are all distinct. In any proposed minor model
in a star, a cross-edge witness for the first edge forces the
center $r$ into one of its two endpoint branch sets. The second
edge likewise forces $r$ into one of its own two branch sets.
These are distinct branch sets, so their disjointness is violated.
This proves actual grid-minor exclusion for every $k\geq2$ and
every star size, without assuming any desired minor count.

\newpage
\section{All finite and asymptotic bounds are ruled out}
The formal result is stronger than a single small-$k$
counterexample. For every real function $F:\mathbb N\to\mathbb R$
and every proposed asymptotic cutoff $K$, choose
\[
 k=\max(K,2),\qquad n\in\mathbb N\text{ with }n>F(k).
\]
The Archimedean property supplies such an actual natural $n$.
The actual star with $n$ leaves is a connected tree, has no
$k\times k$ grid minor, and has edge count $n>F(k)$.
Thus no real-valued function of $k$ alone is an eventual upper
bound, even after restricting to stars. This covers every finite
envelope of the form claimed in $2^{O(k\log k)}$, including
arbitrary constants and cutoffs.

The final Lean theorem explicitly negates a bound over
\emph{all actual finite simple graphs}: such a bound would
restrict to this actual star family and contradict the preceding
theorem. The finite graphs, edge sets, cardinalities, branch
models and $k$-grid objects all appear in the formal statements.
No treewidth value is assigned or needed in the argument.

\section{Primary reference and formal verification}
Chandra Chekuri and Julia Chuzhoy,
\emph{Polynomial Bounds for the Grid-Minor Theorem},
\href{https://arxiv.org/pdf/1305.6577v5}{arXiv:1305.6577v5},
9 August 2016, printed page 4, gives the standard branch-set
model with single cross-edges and the horizontal/vertical grid
definition. The paper states a path-based variant and explicitly
identifies the standard single-edge variant as equivalent.
We adopt the latter standard definition directly; we do not claim
to re-formalize the general equivalence with a sequence of
deletions and contractions. Its abstract and printed pages 1--2
make clear that the real grid theorem concerns treewidth.
The cited passages were directly read by both the worker and
the separate supervising AI.

The package pins Lean 4.33.0 and Mathlib commit
\begin{center}\small\texttt{db584cd6d46c92f209a44c0f1c829460d327499d}\end{center}
and contains the full bilingual original source, complete
LaTeX/PDF, a new source-bound root build, all named declaration
axiom audits and full official
\texttt{leanchecker --fresh Main} replay. The audits explicitly
include the proof-bearing \texttt{MinorModel} structure,
constructor and all branch-condition projections, as well as
every other named local declaration. No \texttt{sorry},
\texttt{native\_decide} or additional axiom is used.

OpenAI Codex assisted, with separate supervising AI source/proof/
report/PDF review. This is not human certification; organizer
acceptance is pending. Only the source's literal finite
edge-count upper bound is refuted, not the genuine treewidth
theorem or the separate lower-bound existence statement.
\end{document}
